\chapter{Kết luận và hướng phát triển}
\label{chap:ketluan}

%==============================================================================
\section{Những việc đã làm được}
\label{sec:da-lam}
%==============================================================================

Báo cáo đặt ra năm nhiệm vụ ở Mục~\ref{sec:muc-tieu}. Phần này đối chiếu từng
nhiệm vụ với kết quả thực tế, kể cả những chỗ chưa đạt.

\paragraph{MT1 --- Cơ sở toán học của mạng truyền thẳng.}
Đã hoàn thành, với trọng tâm khác tài liệu học sâu tổng quát đúng như dự kiến.
Hai hệ thức truy hồi \eqref{eq:G-truy-hoi} và \eqref{eq:S-truy-hoi} cho Jacobi và
Hessian \emph{theo biến đầu vào} được thiết lập tường minh
(Mệnh đề~\ref{md:jacobian}, \ref{md:hessian}) và trở thành công cụ trung tâm cho
mọi tính toán phần dư về sau. Từ chúng suy ra ngay điều kiện \emph{cần} trên hàm
kích hoạt, mà trường hợp ReLU là dạng cực đoan: hàm mục tiêu phần dư của phương
trình bậc hai \emph{hằng số địa phương}, nên mọi thuật toán gradient đứng yên
tuyệt đối (Hệ quả~\ref{hq:pinn-hong}).

\paragraph{MT2 --- Vi phân tự động và thuật toán tối ưu.}
Đã hoàn thành. Mắt xích quan trọng nhất là Mệnh đề~\ref{md:ad-tai-tao}: cơ chế
thuận-trên-thuận của vi phân tự động \emph{tái tạo đúng} hai hệ thức truy hồi của
Chương~\ref{chap:fnn} mà không cần lập trình tường minh $\sigma''$. Nhờ đó,
Thuật toán~\ref{alg:dao-ham} đóng vai trò đặc tả toán học, còn cài đặt thực tế
chỉ là vài dòng gọi \texttt{torch.autograd.grad} với cờ
\texttt{create\_graph=True}. Chiến lược hai giai đoạn Adam~$\to$~L-BFGS được biện
minh bằng ba lập luận cụ thể --- tính tất định của gradient, sàn nhiễu
(Mệnh đề~\ref{md:san-nhieu}), và điều kiện xấu --- chứ không bằng kinh nghiệm
thực hành.

\paragraph{MT3 --- Xây dựng PINNs một cách chặt chẽ.}
Đã hoàn thành, và đây là chỗ báo cáo bổ sung được một mắt xích mà tài liệu
thường bỏ qua: \emph{vì sao cực tiểu hoá phần dư lại là việc đáng làm}.
Mệnh đề~\ref{md:chan-on-dinh} trả lời bằng một chặn ổn định, và
Mệnh đề~\ref{md:chan-suy-bien} chỉ ra hằng số của chặn ấy suy biến theo $1/\nu$.
Đóng góp tổng hợp của chương là Bảng~\ref{tab:ba-nguyen-nhan}: ba nguyên nhân độc
lập, nằm ở ba tầng khác nhau, cùng làm PINNs suy giảm trên đúng lớp bài toán có
thang độ dài nhỏ.

\paragraph{MT4 --- Cài đặt và thực nghiệm trên PyTorch.}
Đã hoàn thành. Toàn bộ cài đặt dùng PyTorch thuần, không dùng thư viện PINNs đóng
gói sẵn (lý do ở Nhận xét~\ref{nx:vi-sao-tu-viet}), và được công bố kèm báo cáo
(Phụ lục~\ref{app:code}) để mọi bảng số và mọi hình trong
Chương~\ref{chap:thucnghiem} đều tái lập được bằng một lệnh.

Phần ``chi phí tính toán thực đo'' mà nhiệm vụ này hứa được trả ở
Mục~\ref{sec:tn9}, và câu trả lời cần được nêu thẳng vì nó bất lợi cho phương
pháp: trên bài toán Burgers một chiều, PINN cần khoảng $170$ giây để đạt sai số
$1{,}38\times10^{-1}$, trong khi bộ giải sai phân hữu hạn đạt cùng mức sai số
trong $0{,}0098$ giây --- chậm hơn khoảng $1{,}8\times10^{4}$ lần, tức bốn bậc độ
lớn. Đây không phải mâu thuẫn với Chương~\ref{chap:modau}: cả ba động lực
nêu ở Mục~\ref{subsec:gioi-han-co-dien} --- bài toán ngược, điều kiện biên
khuyết, lời nguyền số chiều --- đều \emph{không} có mặt trong lớp bài toán mà
báo cáo thực nghiệm. Kết luận đúng phạm vi là: các thí nghiệm ở đây kiểm chứng
\emph{cơ chế} của PINNs, không kiểm chứng \emph{tính cạnh tranh} của chúng
(Nhận xét~\ref{nx:doc-tn9}).

Mục~\ref{sec:tn10} bổ sung vế còn lại của bức tranh: với cấu hình của công trình
gốc và lấy mẫu thích nghi, sai số Burgers trung vị giảm $213$ lần, xuống
$6{,}48\times10^{-4}$ --- xác nhận rằng con số của Mục~\ref{sec:tn5} là giới hạn
của ngân sách chứ không của phương pháp. Nhưng tỉ số chi phí so với sai phân hữu
hạn vẫn ở bậc $10^{4}$ ($1{,}1\times10^{4}$ lần, so với $1{,}8\times10^{4}$ ở
Mục~\ref{sec:tn9}): PINN \emph{có thể} chính xác trên bài toán cứng này, chỉ
không rẻ hơn. Thí nghiệm tách biến đi kèm còn chỉ ra rằng mức cải thiện ấy gần
như hoàn toàn đến từ cấu hình mạng, số điểm phối trí và độ dài L-BFGS. Ba chiến
lược lấy mẫu thích nghi --- thêm điểm, và hai biến thể lấy mẫu lại toàn bộ ---
không chiến lược nào cải thiện nhất quán sai số toàn cục, và cả ba đều làm vùng
trơn tệ đi ở mọi hạt giống: ở cấu hình này, cách đặt điểm phối trí chỉ phân phối
lại sai số giữa các vùng.

\paragraph{MT5 --- Khảo sát độ nhạy và phân tích hạn chế.}
Đã hoàn thành ở mức tương xứng với ngân sách tính toán, nhưng đây cũng là nhiệm
vụ còn nhiều dư địa nhất; Mục~\ref{sec:hanche} đã nêu thẳng chín hạn chế của chính
thiết kế thực nghiệm. Phần khảo sát độ nhạy đi xa hơn dự kiến ở một điểm:
Mục~\ref{sec:tn8} không chỉ đo độ nhạy theo hạt giống mà còn dùng chính phép đo
ấy để \emph{loại bỏ} bốn con số mà các mục trước đã rút ra --- đúng vai trò mà
một khảo sát độ nhạy phải có.

%==============================================================================
\section{Đóng góp của báo cáo}
\label{sec:dong-gop-cuoi}
%==============================================================================

Trong phạm vi đã nêu ở Mục~\ref{subsec:pham-vi}, báo cáo có ba đóng góp.

\emph{Thứ nhất}, một trình bày mạch lạc bằng tiếng Việt về PINNs, đi từ cơ sở
toán học của mạng truyền thẳng tới công thức đầy đủ của mô hình thời gian liên
tục, với notation nhất quán xuyên suốt và mọi khẳng định đều truy được về một
mệnh đề cụ thể.

\emph{Thứ hai}, một cài đặt PyTorch độc lập kèm kết quả thực nghiệm định lượng,
tái lập được, có thể dùng làm điểm khởi đầu cho các nghiên cứu tiếp theo.

\emph{Thứ ba}, một phân tích đối chiếu có hệ thống giữa lý thuyết và số liệu
(Bảng~\ref{tab:doi-chieu}). Phần đối chiếu này cho một kết luận phương pháp luận
đáng chú ý: trong mười bốn hạng mục, \emph{không} một mệnh đề đã chứng minh nào
bị số liệu bác bỏ. Cả bảy hạng mục phải điều chỉnh đều đến từ bốn kiểu suy luận
yếu hơn --- đối chiếu một số mũ đo ở ngân sách hữu hạn với số mũ tiệm cận, ghép
hai kết quả rời rạc, khái quát hoá từ một lần chạy duy nhất, và ước lượng bậc độ
lớn chưa kiểm chứng.

Bên cạnh ba đóng góp đã dự kiến, quá trình thực nghiệm sinh thêm năm kết quả
không nằm trong kế hoạch ban đầu, và cả năm đều xuất phát từ \emph{kết quả tiêu
cực}:

\begin{itemize}[leftmargin=*]
\item \textbf{Một điều chỉnh phát biểu lý thuyết.} Mục~\ref{subsec:he-qua-pinn}
      kết luận rằng bài toán đòi hỏi độ chính xác cao nhất ở đúng những tần số mà
      mô hình học chậm nhất. Thí nghiệm~\ref{tn:pho} cho thấy phát biểu này
      không đúng phổ quát: với toán tử elliptic thuần tuý, hàm mục tiêu phần dư
      \emph{ưu tiên} tần số cao và chính tần số thấp mới bị bỏ rơi. Cách phát
      biểu đã được siết lại ở Nhận xét~\ref{nx:hieu-chinh}.
\item \textbf{Một cách đọc bị siết lại về tầng thứ nhất của
      Bảng~\ref{tab:ba-nguyen-nhan}.} Mục~\ref{sec:tn7} xác nhận chặn của
      Mệnh đề~\ref{md:chan-suy-bien} ở cả $15$ trường hợp, nhưng không tìm được
      bằng chứng rằng tỉ số quan sát được tăng theo $1/\nu$: số mũ ước lượng
      trải từ $-0{,}730$ tới $+0{,}049$ tuỳ hạt giống. Hệ quả:
      hằng số ổn định là một \emph{chặn} đã chứng minh, không phải một
      \emph{cơ chế} đã quan sát (Nhận xét~\ref{nx:doc-tn7}).
\item \textbf{Một ước lượng định lượng của Chương~\ref{chap:autodiff} bị sửa.}
      Nhận xét~\ref{nx:chi-phi-bac-cao} đoán rằng đồ thị dùng cho lượt tối ưu
      dài gấp ``bốn đến sáu lần'' đồ thị xuôi. Mục~\ref{sec:tn9} đo trên mười
      lần chạy độc lập và được trung vị $8{,}4$; hơn nữa nhiễu đo chỉ có thể kéo
      tỉ số \emph{xuống}, nên ước lượng ít chệch nhất là khoảng $9$. Nguyên nhân
      của sai lệch truy được: nhận xét ngầm coi hai bậc đạo hàm tốn như nhau,
      trong khi lượt lấy $\partial_{xx}u$ chạy ngược trên đồ thị \emph{đã được}
      lượt lấy $\partial_{x}u$ nối dài. Điều đáng nói là ước lượng này đã đứng
      trong chương lý thuyết mà không ai kiểm chứng, và chỉ một phép đo thời
      gian đơn giản là đủ để phát hiện.

\item \textbf{Một chế độ hỏng chứng minh được của quy tắc cân bằng trọng số.}
      Thuật toán~\ref{alg:annealing} làm giảm độ chính xác trên cả hai bài toán
      khảo sát. Truy nguyên dẫn tới Bổ đề~\ref{bd:phan-hoi-duong}: quy tắc
      \eqref{eq:lambda-hat} có một vòng phản hồi dương với tốc độ phân kỳ
      $\hat\lambda_i \gtrsim J_i^{-1/2}$ và không có cơ chế tự hãm. Đây là kết
      quả lý thuyết nảy sinh \emph{từ} thực nghiệm chứ không phải được kiểm chứng
      bởi thực nghiệm.
\item \textbf{Bốn con số của chính báo cáo này bị chính nó loại bỏ.}
      Mục~\ref{sec:tn8} lặp lại mọi thí nghiệm có phụ thuộc ngẫu nhiên trên năm
      hạt giống, và kết quả là bốn đại lượng từng được trích dẫn như số đo đều
      hoá ra không tiêu biểu: tương quan $\rho_b$--$\norm{\bW^{(1)}}_F^2$ chỉ còn
      hệ số $0{,}24$ trên thang log (Nhận xét~\ref{nx:tuong-quan-hong}); chiều
      biến thiên của $\rho_b$ trên bài toán lành mạnh chỉ đúng ở $3/5$ hạt giống;
      số mũ $-0{,}907$ của Mục~\ref{sec:tn2} là cực trị của dải
      $[-0{,}907;\ -0{,}503]$; và tỉ số $38{,}6$ của Mục~\ref{sec:tn3} thuộc đuôi
      hiếm của một phân bố có trung vị $5{,}5$. Trong cả bốn trường hợp,
      \emph{kết luận định tính vẫn đứng vững} --- điều bị loại bỏ chỉ là con số.
      Giá trị của kết quả này nằm ở \emph{quy trình}: chi phí của việc lặp năm
      hạt giống là chưa tới hai giờ máy, rẻ hơn hẳn cái giá của bốn con số sai.
\end{itemize}

%==============================================================================
\section{Hạn chế}
\label{sec:han-che-cuoi}
%==============================================================================

Ngoài chín hạn chế về thiết kế thực nghiệm đã nêu ở Mục~\ref{sec:hanche}, cần nói
rõ ba giới hạn về phạm vi của toàn bộ báo cáo.

Thứ nhất, báo cáo \emph{không} chứng minh định lý mới về hội tụ hay ước lượng
sai số cho PINNs. Các kết quả lớn của lĩnh vực chỉ được trích dẫn cùng diễn giải;
những chứng minh được trình bày đầy đủ đều sơ cấp, và giá trị của chúng nằm ở chỗ
biến những khẳng định mơ hồ thành phát biểu có giả thiết rõ ràng.

Thứ hai, phạm vi bài toán hẹp: chỉ mô hình thời gian liên tục, chỉ mạng truyền
thẳng kết nối đầy đủ, chỉ bài toán thuận, và chỉ các phương trình một chiều theo
không gian. Lợi thế lý thuyết quan trọng nhất của mạng nơ-ron --- bậc xấp xỉ
$\mathcal{O}(m^{-1/2})$ không phụ thuộc số chiều ở Định lý~\ref{dl:barron} ---
vì vậy \emph{không} được kiểm chứng bằng thực nghiệm nào trong báo cáo. Phép đo
chi phí ở Mục~\ref{sec:tn9} cho thấy cái giá của giới hạn phạm vi này là cụ thể
chứ không trừu tượng: ở $d = 1$, nơi phương pháp lưới gần như tối ưu, PINN thua
bốn bậc về chi phí, và báo cáo không có bài toán nào để cho thấy tỉ số ấy đảo
chiều.

Thứ ba, các chặn ổn định ở Mệnh đề~\ref{md:chan-on-dinh} và
Mệnh đề~\ref{md:chan-suy-bien} được thiết lập cho bài toán tuyến tính một chiều
với điều kiện biên thoả chính xác. Với ràng buộc mềm và với toán tử phi tuyến,
chúng chỉ đóng vai trò cơ sở khái niệm chứ không phải chặn hậu nghiệm áp dụng
được trực tiếp.

%==============================================================================
\section{Hướng phát triển}
\label{sec:huong-phat-trien}
%==============================================================================

Bảng~\ref{tab:ba-nguyen-nhan} gợi ý một nguyên tắc sắp xếp các hướng tiếp theo:
vì ba nguyên nhân nằm ở ba tầng khác nhau, mỗi hướng nên được đánh giá theo
\emph{tầng mà nó tác động tới}, và không hướng nào thay thế được hai tầng còn
lại.

\paragraph{Tầng bài toán.} Phân rã miền và hành quân theo thời gian
(Bảng~\ref{tab:remedies}) thay một bài toán khó bằng một dãy bài toán dễ hơn, tức
tác động trực tiếp lên hằng số ổn định. Đây là hướng đáng thử trước với phương
trình Burgers độ nhớt nhỏ, nơi Mục~\ref{sec:tn5} cho thấy sai số tập trung vào
một dải hẹp quanh lớp trong.

\paragraph{Tầng tối ưu.} Bổ đề~\ref{bd:phan-hoi-duong} không chỉ ra rằng cân bằng
trọng số là sai, mà chỉ ra \emph{cơ chế} hỏng. Hai sửa đổi đi thẳng từ cơ chế
ấy đáng được khảo sát: đặt chặn trên cho $\lambda_i$, hoặc thay mẫu số của
\eqref{eq:lambda-hat} bằng một đại lượng không triệt tiêu cùng $J_i$ --- chẳng
hạn vết của nhân tiếp tuyến.

\paragraph{Tầng xấp xỉ.} Thí nghiệm~\ref{tn:pho} cho thấy nhúng đặc trưng Fourier
cải thiện $161$ lần nhưng chỉ \emph{dịch chuyển} dải tần hiệu dụng chứ không mở
rộng nó, và việc chọn thang $s$ đòi hỏi biết trước thang độ dài của nghiệm chưa
biết. Một quy tắc chọn $s$ dựa trên dữ liệu --- ví dụ ước lượng từ phổ của phần
dư trong vài trăm vòng lặp đầu --- là hướng cụ thể và khả thi.

\paragraph{Củng cố phương pháp luận thực nghiệm.} Việc đầu tiên trong nhóm này
đã được làm trọn: Mục~\ref{sec:tn8} lặp TN2, TN3, TN4, TN5 trên năm hạt giống,
và riêng nó đã loại được bốn con số không tiêu biểu. Việc còn lại là
\emph{cách trình bày} --- gộp trung vị và khoảng vào chính các bảng của
Mục~\ref{sec:tn1}--\ref{sec:tn7} thay vì để riêng ở Mục~\ref{sec:tn8} --- và
nâng $n$ từ $5$ lên $10$--$20$ cho những đại lượng lệch mạnh. Hai việc rẻ khác:
quét số điểm phối trí $N_r$ cho cả bốn cách xử lý điểm ở
Mục~\ref{subsec:tn10-rad} --- vì kết quả ở đó cho thấy lấy mẫu thích nghi chỉ
phân phối lại sai số ở $N_r = 10\,000$, và câu hỏi còn mở là lợi ích mà tài liệu
báo cáo có xuất hiện lại khi $N_r$ nhỏ hay không;
và tách riêng tầng thứ nhất của Bảng~\ref{tab:ba-nguyen-nhan} bằng một thiết kế
khác với Mục~\ref{sec:tn7}. Phép quét $\nu$ ở đó cho thấy tỉ số
$\norm{e_{\btheta}}/\norm{r_{\btheta}}$ do huấn luyện sinh ra không đo được hằng
số ổn định, vì hướng của sai số do chính quá trình tối ưu quyết định chứ không
phải do bài toán. Cách sạch hơn là \emph{dựng} sai số theo hướng hàm riêng
$\sin \pi x$ --- hướng làm chặn Poincaré trở nên chặt --- rồi đo trực tiếp tỉ số
ấy theo $\nu$.

\paragraph{Mở rộng phạm vi.} Cuối cùng, hai mở rộng nằm ngoài phạm vi báo cáo
nhưng nối liền với nó: bài toán ngược, nơi khung hàm mất mát \eqref{eq:total-loss}
đã sẵn sàng và chỉ cần thêm $\bzeta$ vào tập biến tối ưu
(Mục~\ref{subsec:thuan-nguoc}); và bài toán nhiều chiều, nơi lập luận chống lời
nguyền số chiều ở Nhận xét~\ref{nx:loi-nguyen} lần đầu tiên trở nên kiểm chứng
được. Mở rộng thứ hai là mở rộng \emph{quyết định}, và Mục~\ref{sec:tn9} cho nó
một tiêu chí sắc: lặp lại chính phép đo chi phí ấy khi $d$ tăng, rồi tìm giá trị
$d$ mà tỉ số bốn bậc độ lớn ấy đảo chiều. Chừng nào chưa có con số đó, ưu thế của PINNs
so với phương pháp lưới vẫn là một lập luận tiệm cận chứ chưa phải một quan sát.
